\documentclass[12pt]{report}
\usepackage[francais]{babel}
\usepackage{lmodern}
\usepackage[a4paper]{geometry}
\usepackage[T1]{fontenc}
\usepackage[utf8]{inputenc}  
\usepackage{moreverb}
\usepackage{amsmath}
\usepackage{amsfonts}
\usepackage{amssymb}
\usepackage{textcomp}
\usepackage{pifont}
\usepackage{geometry}
\usepackage[pdftex]{graphicx}
\usepackage{graphics}
\usepackage{url}
\usepackage{graphicx}
\usepackage{float}
\usepackage{color}
\usepackage[nottoc, notlof, notlot]{tocbibind}
\usepackage[french]{varioref}
\usepackage[Glenn]{fncychap}
\usepackage{pdfpages}

\usepackage{multicol}

% entête et pied de page
%% \usepackage{fancyhdr} 
%% \pagestyle{fancy}
%% \renewcommand{\chaptermark}[1]{\markboth{#1}{}}
%% \renewcommand{\sectionmark}[1]{\markright{\thesection\ #1}}
%% \fancyhf{} \fancyhead[LE,RO]{\bfseries\thepage}
%% \fancyhead[LO]{\bfseries\rightmark}
%% \fancyhead[RE]{\bfseries\leftmark}
%% \renewcommand{\headrulewidth}{0.5pt}
%% \addtolength{\headheight}{0.5pt}
%% \renewcommand{\footrulewidth}{0pt}
%% \fancypagestyle{plain}{ \fancyhead{}
%% \renewcommand{\headrulewidth}{0pt}} 

% pour inclure du code par exemple
\usepackage{listings}

\lstset{%configuration de listings
float=hbp,%
basicstyle=\ttfamily\small, %
columns=flexible, %
tabsize=2, %
frame=trBL, %
frameround=tttt, %
extendedchars=true, %
showspaces=false, %
showstringspaces=false, %
numbers=left, %
numberstyle=\tiny, %
breaklines=true, %
breakautoindent=true, %
captionpos=b,%
xrightmargin=0cm, %
xleftmargin=-0cm, %
language=tex, %
frameround=fttt;%
}

%%%%%%%%%%%%%%%% Lengths %%%%%%%%%%%%%%%%
\geometry{a4paper,twoside,left=2cm,right=2cm,marginparwidth=1.2cm,marginparsep=3mm,top=1.7cm,bottom=1.5cm}

\newcommand{\stamp}{{\tt \textit{Stamp }}}
\newcommand{\class}{{\tt \textit{class }}}
\newcommand{\initarg}{{\tt \textit{Initarg }}}
\bibliographystyle{plain}
\urlstyle{sf}

%%%%%%%%%%%%%%%%%%%%%%%%%%%%%%%%%%%%%%%%%%%%%%%

\newenvironment{vcenterpage}
{\newpage\vspace*{\fill}}
{\vspace*{\fill}\par\pagebreak}

\newtheorem{ex}{Exemple}%[section]
\newtheorem{theo}{Theorem}
\newcommand{\tuple}[1]{\ensuremath{\langle #1 \rangle}}

\begin{document}
%%%% Page de titre %%%%
\def\logo{
  \begin {figure}[H]
	\includegraphics[scale=0.33]{\DIR/img/logo.jpg}
        \hspace{1cm}
        \includegraphics[scale=0.2]{\DIR/img/logobordeaux1.jpg}
        \hspace{2.8cm}
        \includegraphics[scale=1]{\DIR/img/logoEvollis.png}        
        \hspace{0.8cm}   
        \includegraphics[scale=0.3]{\DIR/img/logoLaBRI.jpeg}	

	\label{logo}
  \end {figure}
}

\def\title{Mise en place d'une solution non relationnelle de gestion de données}

\def\intervenant{
        \begin{flushleft}
	  \begin{tabbing}
		\textbf{Maître de stage:}
                \hspace{7.2cm} \=\textbf{Étudiant:} \\
                \noindent Eric de {\sc Marignan}
                \> {\sc Barro} Lissy Maxime\\
                \> \\
                \noindent \textbf{Tuteur à} {\sc bordeaux}1 {\bf :}
                \> Élève ingénieur\\
                \noindent Mohamed {\sc Mosbah}
                \> Dernière année - Option GL \\ 
                \noindent Sofian {\sc Maabout} \> {\sc enseirb-matmeca} 2011/2012 \\
                \>\\
                \noindent \textbf{Tuteur à l'{\sc enseirb}:}
                \> \\
                %\noindent aaaa {\sc aaaa}
	  \end{tabbing}
        \end{flushleft}
}

\def\date{\today}%9 juin 201}

\def\job{PROJET DE FIN D'ÉTUDE / STAGE MASTER 2 RECHERCHE}

\begin{titlepage}
  \logo
  \begin{flushleft}
    \textbf{École Nationale Supérieure d’Électronique, Informatique,
      Télécommunications, Mathématique et Mécanique de Bordeaux}

    \vspace{0.5cm}

    \textsf{Département informatique}

    \vspace{0.5cm}

   1 avenue du Dr Albert Schweitzer\\
   B.P. 99 33402 Talence Cedex

    
  \end{flushleft}
  
  \vspace{4cm}
	\begin{center}
	  {\bf \job}\\
	  \vspace{1cm}
		 {\LARGE\bf \title}\\
\vspace{1cm}
           \date

	\end{center}


        \vspace{4cm}
        \intervenant
\end{titlepage}


%%%%% terminologie %%%%%%%%%%
\newpage

\def\termea{\bf \footnotesize xx}
\def\sensa{\footnotesize xx}

\def\termeb{\bf \footnotesize xx}
\def\sensb{\footnotesize xx}

\def\termec{\bf \footnotesize xx}
\def\sensc{\footnotesize xx}

\def\termed{\bf \footnotesize }
\def\sensd{\footnotesize }

\def\termee{\bf \footnotesize }
\def\sense{\footnotesize }

\def\termef{\bf \footnotesize }
\def\sensf{\footnotesize }

\def\termeg{\bf \footnotesize }
\def\sensg{\footnotesize }

\def\termeh{\bf \footnotesize }
\def\sensh{\footnotesize }

\def\termei{\bf \footnotesize }
\def\sensi{\footnotesize }

\def\termej{\bf \footnotesize }
\def\sensj{\footnotesize }

\def\termek{\bf \footnotesize }
\def\sensk{\footnotesize }



\begin{center}
\subsubsection*{Terminologie}
\begin{tabular}{|p{5cm}|p{12cm}|}
\hline
{\bf ~~~~~ T{\scriptsize ERME}} &  {\bf ~~~~~~~~~~~~~~~~~~~ S{\scriptsize IGNIFICATION}}\\
\hline
\hline
\termea & \sensa\\
\hline
\termeb & \sensb\\
\hline
\termec & \sensc\\
\hline
\termed & \sensd\\
\hline
\termee & \sense\\
\hline
\termef & \sensf\\
\hline
\termeg & \sensg\\
\hline
\termeh & \sensh\\
\hline
\termei & \sensi\\
\hline
\termej & \sensj\\
\hline
\termek & \sensk\\
\hline

\end{tabular}
 
\end{center}


%%%% resume %%%%%%%%%
\begin{abstract}
  %% français

%% \begin{center}
%% \subsubsection*{Abstract}
%% \end{center}

%% \noindent english 

\end{abstract}

%%%% plan %%%%%
\tableofcontents
\listoffigures
\listoftables

%%%% remerciement %%%%%%%
\addcontentsline{toc}{chapter}{Remerciements}
\chapter*{Thanks}


%%%%% corps du rapport %%%%%%%%%
\chapter*{Introduction}

The article ~\cite{pd} talks about how to visualize graphs containing many nodes and edges. Improvements in data acquisition leads to an increase of the size and the complexity of graphs and this huge amount of data generally causes visual clutter, in our case due to edge crossing.
For example, it could be interesting to visualize data in fields like biology, social sciences, data mining or computer science, and then emphasize their high-level pattern to help users perceive underlying models.


Nowadays, in the research world, the information is easily represented into graphs to visualize more and more data. However, this huge amount of information prevents the graph from being manually drawn:  It explains the need of automatic methods able to generate an appropriate graph with all nodes and edges. Yet this graph may suffer from cluttering, which should be reduced for a better understanding.

Our objective all along this project is to read what has been done before relating to this problem, to provide an objective point of view on those previous works, and propose our contribution. We have implemented a method, then optimized its performances with
 current technologies (OpenMP, Tulip...) and setted our boundaries. 


The first part of this document presents review-related work on reducing edge clutters and enhancing edge bundle visualization, with which the article is connected. The second will deal with the Tutte algorithm and its differents versions. A third part will talk about the implementation issues and show our results. Finally, we draw a conclusion and explain the limits of our work for further improvements.

\addcontentsline{toc}{chapter}{Introduction}

\part{Présentation du cadre de mon stage}

\chapter{Présentation de l'organisme d'acceuil: {\bf EVOLLIS}}


\chapter{Présentation du projet XXXXX}
\section{Le contexte}
\section{L'environnement de travail}
\section{L'équipe de travail}

\part{Déroulement de mon stage}
\chapter{État de l'art}

\chapter{Le noSQL mongoDB}

\chapter{Le noSQL de LaBRI}

\chapter{Les tests de charge}

\chapter{Pilote m ou c pour EVOLLIS}

\addcontentsline{toc}{chapter}{Conclusion}
% 


\begin{thebibliography}{}

\bibitem{pd} A. Lambert, R. Bourqui, and D. Auber. Winding roads: Routing edges
into bundles. \emph{In 12th Eurographics/IEEE-VGTC Symposium on
Visualization (Computer Graphics Forum; Proceedings of EuroVis
2009).} To appear., 2010.

\bibitem{pe} A. Lambert, R. Bourqui, and D. Auber. 3D edge bundling for
geographical data visualization. \emph{In IV ’10: Proceedings of the 14
International Conference on Information Visualisation (IV’09)},
Washington, DC, USA, 2010. IEEE Computer Society.

\bibitem{pb} B. Bollob's. Modern graph theory, \emph{volume 184 of Graduate Texts in Mathematics}. Springer-Verlag, 1998.

\bibitem{pa} E. Colin de Verdière, M. Pocchiola, and G. Vegter. Tutte's Barycenter Method applied to Isotopies. \emph{Computational Geometry: Theory and Applications, 26}, 81–97, 2003.

\bibitem{pf} Gormen, T.H. and Leiserson, C.E. and Rivest, R.L. and Stein, C. Introduction to algorithms. \emph{In MIT press Cambridge, MA}, 16:"Greedy Algorithms", 1990.

\bibitem{pc} William T. Tutte. How to draw a graph. \emph{Proceedings of the London Mathematical Society}, 13:743–768, 1963.

% \bibitem{1} 

% \bibitem{2} 

% \bibitem{3} 

% \bibitem{4} 

% \bibitem{5} 
   
\end{thebibliography}


%%%%%% Liste des annexes %%%%%%%%%%
\part{Annexes}
\appendix
%% \chapter{Les propriétés \textsf{ACID}}\label{acid}
Les propriétés \textsf{ACID} permettent à un \textsf{SGBD} d'effectuer
des transactions. Par transaction, il faut comprendre une suite
d'opérations qui font passer la \textsf{BDD} d'un état antérieur à un
état postérieur. Les états intermédiaires entre les états avant la
transaction et après la transaction ne sont pas visibles. Ses
propriétés sont les suivantes:

\paragraph{Atomicité:} la suite d'opérations constituant une transaction 
est indivisible. La transaction est entièrement effectuée ou pas du tout. 
Il y a annulation de toute la transaction lorsqu'une des opérations échoue.
S'il est question de modifier une série de valeurs et qu'une modification
échoue alors toutes les valeurs déjà modifiées reprennent leurs anciennes
valeurs. 

\paragraph{Cohérence:} quelque soit l'opération effectuée, la base doit garder
un état cohérent. Toute transaction qui viole par exemple une règle d'intégrité 
échoue. Après la fusion de deux tables, les entrées doivent toutes 
avoir des identités différentes. Si ce n'est pas le cas alors la fusion n'est 
pas effectuée.

\paragraph{Isolation:} chaque transaction est isolée de sorte à ce qu'elle est 
seule peut voir les modifications pendant son exécution. Toute
transaction enclenchée en parallèle d'une autre voit la version des
données antérieure à celle-ci.  Il existe 4 niveaux d'isolation
définis dans le standard \textsf{ANSI/ISO SQL}:
\begin{enumerate}
\item \textsf{Uncommited read ou lectures des données non validées}.
  Ce niveau est le niveau d'isolation le plus léger. Avec un tel
  niveau d'isolation le système se comporte comme s'il n'y en avait
  pas. Les modifications apportées par une transaction non validée
  sont visibles.

\item \textsf{Commited read ou lecture des seules données
  validées}. Les modifications lors d'une transaction ne sont visibles
  que lorsqu'elle termine.  Cependant lors d'une transaction une
  donnée peut changer sans que la transaction en cours en soit
  responsable. Ceci peut arriver, par exemple, lorsqu'une transaction
  en parallèle termine et que ses modifications sont validées.

\item \textsf{Repeatable read ou lecture répétée}. Ce niveau
  fonctionne comme le niveau précédent à la seule différence que
  durant son exécution, une transaction ne voit pas les mises à jour
  effectuées par d'éventuelles transactions qui se sont exécutées en
  parallèle. D'où « \textsf{repeatable read} » pour souligner que pendant une
  transaction, une donnée aura toujours la même valeur en
  lecture si elle n'est pas modifiée par la transaction elle-même.
  Cependant tout nouveau rajout validé de données au système par une
  transaction qui a terminé est visible par toute autre transaction en
  cours.

\item \textsf{Serializable ou sérialisable}. Ce niveau est le niveau
  d'isolation le plus poussé. Le système se comporte comme s'il n'y
  avait qu'une transaction à la fois. Pendant son exécution une
  transaction ne voit ni les mises à jour, ni les rajouts de données
  au système des autres transactions. D'où « \textsf{serializable} » pour
  mettre en relief le caractère d'exécution en série des transactions
  plutôt qu'en parallèle.
\end{enumerate}

\paragraph{Durabilité:} dès lors qu'une transaction est validée, aucune défaillance
du système ne pourra conduire à l'annulation de celle-ci. Les modifications liées à
une transaction validée perdurent et ne sont jamais remises en cause.

\end{document}
